% Source: User reference pp. 98-102; framework: https://github.com/tianyicui/pack (V2.lyx).
\subsection{泛化物品与背包模型归纳}
\label{dp:generalized-knapsack}
\subsubsection*{原理}

把一个物品、一个组或一个子树的可行选择，表示成“精确费用 $t$ 下能取得的最大价值”数组 $h_t$，不可达处为 $-INF$。两个独立部分合并时枚举费用分配：
\[
 (f\star h)_j=\max_{0\le t\le j}(f_{j-t}+h_t).
\]
0/1 物品只有费用 0 与 $w$ 两个可选状态；有限物品有 $0,w,\ldots,cw$；组内互斥的多个方案先在相同费用上取最大，再形成一个费用表。树上依赖先保证每个子树表中的选择合法，才能与其他子树组合。

不要因为把一个部分称为“一件泛化物品”，就把合并复杂度当成 $O(V)$；一般两张费用表的合并为 $O(V^2)$。特殊结构、稀疏状态或单位费用选点数可能进一步优化，须单独证明。

实际选型时依次检查：物品可取次数、是否互斥、是否依赖、费用维数、是否恰好占满、目标为最值还是计数、是否需要还原方案。以上各项变化可能同时出现，必须组合它们的状态定义，而非只套同一个循环。
